\scanpage{215}
\margin{\boxed{\text{p.d. dans les anneaux alternés}}}
\margin{Feuille \unsure{jaunie|paumée}...}
\uline{Théorème} Soient $S$ un anneau muni d'un idéal $I$ à puiss. divisées, $M^*=M^0+M^1$ un $S$-module gradué mod $2$,
\[
A(\unsure{M^*})=\Gamma^*(M^0)\otimes_S\Lambda^*M^1
\]
\inserted{$S$-algèbre} graduée par la parité du degré total, [et augmentée vers $S$ de la façon évidente $\varepsilon:A\longrightarrow S$] $J$ \unsure{l'idéal} engendré dans $A$ par $I$ $(\subset S\subset A)$ et par $M^*(\subset\unsure{S|A})$, [donc $J=\varepsilon^{-1}(I)$].

a) On peut trouver sur \deleted{$A$} $J\subset A$ une structure à p.d. telle que $1^\circ)$ $S\longrightarrow A$ soit un hom. d'anneaux compatible aux p.d. sur $I\subset S$ et $J\subset A$ $2^\circ)$ Pour $(S,M^*)$ variable, cette structure à p.d. est fonctorielle. Ces conditions caractérisent de façon unique la structure à p.d. envisagée.

b) Soit $(B^*\deleted{\unsure{K^*}})$ une $S$-algèbre graduée mod $2$, \deleted{\unsure{alternée strict}} alternée, munie d'un idéal $K^*$ gradué à p.d., \deleted{l'application de restriction} \inserted{compatibles avec les p.d.} sur $I\subset S$.
L'application de restriction
\[
\operatorname{Hom}_{S\text{-alg à p.d.}}((A(M^*),J(M^*)),(B^*,K^*))
\longrightarrow\operatorname{Hom}_{S\text{-mod grad}}(M^*,K^*)
\]
est bijective. (Donc il y a lieu d'appeler $A(M^*)$ la $S$-algèbre \inserted{graduée alternée} à p.d. \unsure{sur} $S$ \uline{engendrée} par le $S$-module \uline{gradué} $M^*$.)

\uline{Dém} :
\margin{214}
